\clearpage

%%%%%%%%%%%%%%%%%%%%%%%%%%%%%%%%%%%%%%%%%%%%%%%%%%%%%%%%%%%%%%%%%%%%%%%%%%%%%%%%%
%%        %%%        %%%        %%%        %%%        %%%         %%%  %%%%  %%%
%%  %%%%%%%%%  %%%%%%%%%  %%%%%%%%%%%%  %%%%%%%%%  %%%%%%  %%%%%  %%%    %%  %%%
%%        %%%        %%%  %%%%%%%%%%%%  %%%%%%%%%  %%%%%%  %%%%%  %%%  %  %  %%%
%%%%%%%%  %%%  %%%%%%%%%  %%%%%%%%%%%%  %%%%%%%%%  %%%%%%  %%%%%  %%%  %%    %%%
%%        %%%        %%%        %%%%%%  %%%%%%        %%%         %%%  %%%   %%%
%%%%%%%%%%%%%%%%%%%%%%%%%%%%%%%%%%%%%%%%%%%%%%%%%%%%%%%%%%%%%%%%%%%%%%%%%%%%%%%%%
 \section{\ROOT をコマンドラインで使う}


  \subsection{\ROOT を起動する}
\begin{verbatim}
	$ root
\end{verbatim}
\ROOT の起動画面が立ち上がり、
\begin{figure}[htbp]
 \begin{center}
  \includegraphics[width = 80mm]{./picture/root_starting.eps}
 \end{center}
 \caption{\ROOT の起動画面}
 \label{Fig:root_starting}
\end{figure}
\begin{verbatim}
	*******************************************
	*                                         *
	*        W E L C O M E  to  R O O T       *
	*                                         *
	*   Version   5.34/09      26 June 2013   *
	*                                         *
	*  You are welcome to visit our Web site  *
	*          http://root.cern.ch            *
	*                                         *
	*******************************************

	ROOT 5.34/09 (v5-34-09@v5-34-09, Jun 26 2013, 17:10:36 on macosx64)

	CINT/ROOT C/C++ Interpreter version 5.18.00, July 2, 2010
	Type ? for help. Commands must be C++ statements.
	Enclose multiple statements between { }.
\end{verbatim}
という一連の情報が表示された後、
\begin{verbatim}
	root[0]
\end{verbatim}
となって\ROOT が立ち上がり操作可能になる。

   \subsubsection{\ROOT の起動画面を省略する}
   \verb|root|起動時に\verb|-l|オプションをつければ起動画面は省略される。
\begin{verbatim}
	$ root -l
\end{verbatim}

   \subsubsection{\ROOT の起動画面を常に省略する}
   いちいち \verb|-l| オプションをつけるのが面倒くさい人はbash起動時に自動的に読み込まれるファイル \verb|~/.bash_profile| に次の文章を追加する。
\begin{verbatim}
	alias root="root -l"
\end{verbatim}
次に\verb|source|コマンドを使用すればよい。
\begin{verbatim}
	$ source ~/.bash_profile
\end{verbatim}


この状態でも\ROOT の起動画面を省略しない場合には
\begin{verbatim}
	$ \root 
\end{verbatim}

   \subsubsection{\ROOT がデフォルトで読み込むファイル}

  \subsection{\ROOT を終了する}

\begin{verbatim}
	root[] .q
\end{verbatim}
